% ============================================================
% RESENHA CRITICA - CAPITULO 7: "CIENCIA E IDEOLOGIA"
% Disciplina: Filosofia e Metodologia da Ciencia - UNIFAL-MG
% Discente: Jeann Victor Batista
% Docente: Jose Francisco Lopes Xarao
%
% Estrutura ajustada ao modelo fornecido e as orientacoes da tarefa.
% Versao completa, estruturada conforme o modelo e os criterios da tarefa:
% identificacao, referencia da obra, resumo e palavras-chave, abstract e
% keywords, apresentacao do autor, tema/problema/tese, sintese, apreciacao
% critica, indicacao do publico-alvo e referencias.
% ============================================================

\documentclass[12pt,a4paper]{article}

% ---------- Codificacao, idioma e tipografia ----------
\usepackage[T1]{fontenc}
\usepackage[utf8]{inputenc}
\usepackage[brazilian]{babel}
\usepackage{lmodern}
\usepackage{microtype}

% ---------- Configuracao de pagina (padrao academico/ABNT) ----------
\usepackage[
  top=3cm,
  left=3cm,
  bottom=2cm,
  right=2cm
]{geometry}

% ---------- Paragrafos e espacamento ----------
\usepackage{indentfirst}
\usepackage{setspace}
\setlength{\parindent}{1.25cm}
\setlength{\parskip}{0pt}
\setlength{\emergencystretch}{3em}

% ---------- Recursos adicionais ----------
\usepackage{url}
\usepackage[hidelinks]{hyperref}

\hypersetup{
  pdftitle={Resenha critica do capitulo Ciencia e ideologia},
  pdfauthor={Jeann Victor Batista},
  pdfsubject={Filosofia e Metodologia da Ciencia}
}

% ---------- Comandos auxiliares ----------
\newcommand{\tituloLivro}{A construção das ciências: introdução à filosofia e à ética das ciências}
\newcommand{\tituloCapitulo}{Ciência e ideologia}

\begin{document}

% ============================================================
% IDENTIFICACAO
% ============================================================
\begin{singlespace}
\begin{center}
  \textbf{UNIVERSIDADE FEDERAL DE ALFENAS -- UNIFAL-MG}\\
  Curso de Ciência da Computação\\
  Disciplina: Filosofia e Metodologia da Ciência\\
  Docente: José Francisco Lopes Xarão\\
  Discente: Jeann Victor Batista

  \vspace{0.8cm}

  {\bfseries\large RESENHA CRÍTICA DO CAPÍTULO ``CIÊNCIA E IDEOLOGIA''}
\end{center}

\vspace{0.5cm}

\noindent\textbf{Resenha da obra:} FOUREZ, Gérard. Ciência e ideologia.
\textbf{In}: FOUREZ, Gérard. \textbf{A construção das ciências: introdução à
filosofia e à ética das ciências}. Tradução de Luiz Paulo Rouanet. São Paulo:
Editora da Universidade Estadual Paulista, 1995. cap. 7, p. 179--193.
\end{singlespace}

\vspace{0.6cm}

% ============================================================
% RESUMO TECNICO-CIENTIFICO DA RESENHA
% ============================================================
\begin{singlespace}
\noindent\textbf{Resumo}

Esta resenha analisa o capítulo ``Ciência e ideologia'', de Gérard Fourez,
com o objetivo de apresentar sua concepção dos discursos ideológicos e a
relação contraditória que eles mantêm com a ciência. O capítulo investiga
como o conhecimento científico pode revelar os critérios e interesses
ocultos por uma ideologia, mas também como a própria ciência produz efeitos
ideológicos quando apresenta conceitos históricos e parciais como verdades
neutras e universais. Fourez distingue as ideologias de primeiro e segundo
graus, discute os limites da ciência diante de questões éticas e examina a
influência dos grupos sociais sobre a produção do conhecimento. A apreciação
crítica relaciona essas ideias ao filme \textit{Oppenheimer}, de Christopher
Nolan, especialmente ao conflito entre a possibilidade técnica de construir
a bomba atômica e as decisões políticas e morais relativas ao seu uso.
Conclui-se que a ciência é indispensável para examinar aspectos específicos
da realidade, mas não oferece, isoladamente, respostas completas para
problemas humanos, sociais e éticos.

\vspace{0.3cm}

\noindent\textbf{Palavras-chave:} ciência; ideologia; ética; discurso
científico; construção do conhecimento.

\vspace{0.5cm}

\noindent\textbf{Abstract}

This review analyzes the chapter ``Science and Ideology'', by Gérard Fourez,
with the aim of presenting his conception of ideological discourses and the
contradictory relationship they maintain with science. The chapter
investigates how scientific knowledge can reveal the criteria and interests
concealed by an ideology, as well as how science itself produces ideological
effects when it presents historical and partial concepts as neutral and
universal truths. Fourez distinguishes between first- and second-degree
ideologies, discusses the limits of science in ethical questions, and
examines the influence of social groups on the production of knowledge. The
critical assessment relates these ideas to Christopher Nolan's film
\textit{Oppenheimer}, especially to the conflict between the technical
possibility of building the atomic bomb and the political and moral decisions
concerning its use. It concludes that science is indispensable for examining
specific aspects of reality, but, by itself, it does not provide complete
answers to human, social, and ethical problems.

\vspace{0.3cm}

\noindent\textbf{Keywords:} science; ideology; ethics; scientific discourse;
construction of knowledge.
\end{singlespace}

\vspace{0.7cm}

% ============================================================
% CORPO DA RESENHA
% As secoes seguem a organizacao utilizada na resenha anterior.
% Na versao final, cada secao deve manter uma redacao fluida e coerente.
% ============================================================
\onehalfspacing

% ---------------------------------------------------------
\section*{Apresentação da obra e do autor}

O capítulo \textit{Ciência e ideologia} é de autoria de Gérard Fourez e
integra o livro \textit{\tituloLivro}, publicado originalmente em francês e
traduzido para o português por Luiz Paulo Rouanet. A edição brasileira
utilizada nesta resenha foi publicada pela Editora da Universidade Estadual
Paulista em 1995.

Gérard Fourez foi um intelectual belga com formação em Filosofia e
Matemática e doutorado em Física Teórica. O autor atuou como professor na
Universidade de Namur, onde fundou um departamento voltado ao estudo das
relações entre ciência, filosofia e sociedade. Dedicou parte importante de
seu trabalho à filosofia e à ética das ciências, à epistemologia e ao ensino
científico (EDITORA UNESP, [s.d.]). Essa trajetória interdisciplinar ajuda a
compreender a perspectiva adotada no capítulo, no qual a ciência não é
apresentada como uma atividade neutra e isolada, mas como uma construção
histórica e social ligada a interesses, valores e diferentes visões de mundo.

% ---------------------------------------------------------
\section*{Tema, problema e posição defendida}

O tema central do capítulo é a relação entre ciência e ideologia. Fourez
analisa tanto a capacidade do conhecimento científico de questionar os
discursos ideológicos quanto a possibilidade de a própria ciência funcionar
como uma ideologia. Essa discussão envolve o modo como os conceitos
científicos são construídos, os limites de sua aplicação e a influência dos
contextos históricos e sociais sobre a produção do conhecimento.

O problema que conduz o capítulo pode ser formulado da seguinte maneira: se
a ciência também é uma construção histórica, orientada por paradigmas,
interesses e escolhas, até que ponto ela pode criticar as ideologias sem se
transformar em um discurso ideológico? A essa questão acrescenta-se outra:
os conhecimentos científicos são suficientes para oferecer respostas
neutras e definitivas a problemas éticos e sociais?

A posição defendida por Fourez é a de que a ciência constitui uma ferramenta
poderosa para a crítica das ideologias, pois permite transformar afirmações
gerais em questões específicas que podem ser examinadas e testadas. Contudo,
ela não produz uma visão completa e absolutamente neutra da realidade. Todo
discurso científico conserva um caráter ideológico de primeiro grau quando
assume que seus conceitos são históricos, parciais e construídos. Ele se
torna uma ideologia de segundo grau quando esconde essas condições e
apresenta uma explicação particular como se fosse uma verdade natural,
universal e definitiva. Embora os interesses sociais influenciem a ciência,
isso não significa que qualquer explicação possa ser aceita, pois a
realidade e a experiência impõem limites às construções científicas
(FOUREZ, 1995, p. 179--193).

% ---------------------------------------------------------
\section*{Síntese do conteúdo}

% A sintese acompanha as sete partes do capitulo e evidencia a continuidade
% do argumento: definicao da ideologia, alcance da critica cientifica,
% limites dessa critica e aplicacao do problema a propria ciencia.
O capítulo 7, intitulado \textit{Ciência e ideologia}, está organizado em
sete partes que desenvolvem um mesmo raciocínio. Fourez começa definindo os
discursos ideológicos, passa a examinar como a ciência pode criticá-los e,
em seguida, demonstra os limites dessa crítica. A partir desses limites, o
autor volta sua análise para a própria ciência, mostrando como ela também
pode produzir efeitos ideológicos e ser influenciada pelos grupos sociais.

Na primeira parte, ``Discursos ideológicos e eficácia crítica da ciência'',
o autor define os discursos ideológicos como aqueles que se apresentam como
descrições da realidade, mas que atuam principalmente para legitimar
práticas, motivar pessoas e reforçar a coesão de certos grupos. Para isso,
destacam alguns aspectos e ocultam outros, inclusive os critérios da visão
que apresentam. Fourez exemplifica esse efeito com uma prática na qual
alguém permite que outra pessoa disponha de parte de sua intimidade e
criatividade em troca de dinheiro. Essa descrição costuma ser associada à
prostituição, embora também possa caracterizar um aspecto do contrato de
trabalho. Desse modo, a ideologia não precisa produzir uma afirmação
inteiramente falsa; ela também atua sobre a maneira como as semelhanças e as
diferenças são percebidas (FOUREZ, 1995, p. 179--180).

Depois de apresentar esse funcionamento, Fourez analisa, na segunda parte,
``Crítica da ideologia pela ciência'', de que modo um discurso ideológico
pode ser questionado. A crítica ocorre quando são revelados os pontos de
vista, os interesses e os critérios que estavam escondidos. A ciência
contribui para isso ao transformar conceitos amplos, como inteligência ou
exploração, em definições específicas que possam ser testadas. Esses testes
permitem verificar aspectos concretos e revelar os limites de certas
generalizações, mas não demonstram uma verdade universal sobre o conceito
original. Eles constituem apenas uma tradução parcial do problema. Além
disso, a passagem para uma linguagem técnica pode apagar a experiência ou o
sentimento que originou o discurso, como ocorre quando a insatisfação dos
estudantes com os horários das provas é reduzida a critérios de organização.
Por essa razão, a ciência é uma importante forma de crítica, mas valores e
ideias filosóficas, éticas ou religiosas também podem levar ao
questionamento de uma ideologia (FOUREZ, 1995, p. 180--184).

Na terceira parte, ``Incapacidade da ciência em esclarecer inteiramente as
questões éticas'', o autor aprofunda o limite apresentado anteriormente. Se
os testes científicos respondem apenas a questões delimitadas, eles não
podem decidir sozinhos problemas que dependem de valores e escolhas.
Diferentes áreas podem oferecer dados e definições, mas nenhuma determina,
de maneira absoluta, quem deve receber direitos ou qual decisão ética deve
ser tomada. Essa limitação aparece na diferença entre tradução e redução. A
tradução representa uma experiência ampla por meio de um conceito científico
específico; a redução trata essa representação parcial como se ela explicasse
toda a experiência. Dizer que o amor é apenas uma questão de hormônios, por
exemplo, transforma uma explicação biológica limitada em uma definição
completa de uma experiência humana. Nesse momento, o discurso científico
passa a ocultar outros aspectos do problema e se aproxima do funcionamento
ideológico (FOUREZ, 1995, p. 184--186).

A quarta parte, ``Dois graus de véus ideológicos'', desenvolve essa
aproximação ao mostrar que até os conceitos científicos permanecem ligados
a contextos, interesses e paradigmas. O conceito de desenvolvimento, por
exemplo, conduz a estudos diferentes quando é definido como crescimento
econômico, realização individual ou autonomia das populações mais pobres.
A escolha inicial já privilegia determinado ponto de vista. Fourez chama de
ideologia de primeiro grau o discurso que reconhece a construção histórica,
os critérios e os limites de seus conceitos. Na ideologia de segundo grau,
esses vestígios são apagados e uma visão particular é apresentada como
natural e universal. Na quinta parte, ``A ciência como ideologia'', essa
distinção é aplicada ao conhecimento científico: a ciência permanece no
primeiro grau quando admite seus limites e passa ao segundo quando se
apresenta como totalmente neutra e definitiva. É a diferença entre afirmar
que o carro pode ser uma resposta a um problema de transporte e tratá-lo
como a única resposta possível. Assim, a ciência deixa de exercer sua função
crítica quando transforma uma conclusão particular em uma explicação total
da realidade (FOUREZ, 1995, p. 186--189).

Na sexta parte, ``O caráter não consciente e implícito das ideologias e a
ética diante das ideologias'', Fourez mostra que esse processo geralmente
não é intencional. Uma pessoa pode apresentar uma definição biológica da
vida como se fosse completa ou reproduzir estereótipos mesmo acreditando
possuir outra visão. A propaganda é diferente porque nela a ocultação é
utilizada conscientemente para atender a projetos políticos ou econômicos.
Como toda representação do mundo é influenciada por critérios e pelo meio
social, o autor considera impossível eliminar completamente as ideologias.
A questão ética consiste em analisar aquilo que transmitimos e decidir quais
ideologias estamos dispostos a assumir ou rejeitar. Essa responsabilidade é
especialmente importante para cientistas e professores, cujas atividades
também comunicam uma visão de mundo (FOUREZ, 1995, p. 189--191).

Na sétima e última parte, ``A ciência varia de acordo com o grupo social?'',
Fourez amplia a discussão para a origem social dos conhecimentos
científicos. A ciência moderna surgiu em uma civilização específica e foi
influenciada pela visão de mundo da burguesia, especialmente pelo projeto de
conhecer, dominar e explorar a natureza. Essa influência também aparece nas
prioridades das disciplinas: a medicina, por exemplo, teria dado maior
atenção ao tratamento individual das doenças do que à prevenção e à higiene
por ter se estruturado em torno das necessidades das classes privilegiadas.
Outros grupos podem desenvolver saberes e prioridades diferentes, mas isso
não significa que cada grupo possa inventar qualquer ciência conforme seus
interesses. Os projetos sociais orientam as perguntas e as aplicações, mas a
realidade impõe resistências e pode demonstrar o fracasso de certas
construções. O capítulo conclui, portanto, que a ciência é uma construção
histórica e social, sem ser por isso uma invenção arbitrária (FOUREZ, 1995,
p. 191--193).

% ---------------------------------------------------------
\section*{Apreciação crítica}

% A apreciacao relaciona o capitulo a uma obra conhecida, avalia o alcance
% de seus exemplos e conclui com a indicacao justificada do publico-alvo.
A leitura do capítulo remete ao filme \textit{Oppenheimer}, de Christopher
Nolan, que apresenta, entre outros acontecimentos, a criação da bomba
atômica durante o Projeto Manhattan. Uma das justificativas iniciais para o
desenvolvimento da arma era o receio de que a Alemanha nazista conseguisse
produzi-la primeiro. Entretanto, a Alemanha foi derrotada antes da conclusão
do projeto e sem ter construído a bomba. Ainda assim, o trabalho prosseguiu,
evidenciando que aquele conhecimento científico não estava relacionado
apenas à ameaça alemã, mas também aos interesses militares e políticos dos
Estados Unidos (OPPENHEIMER, 2023).

Os físicos envolvidos possuíam os conhecimentos necessários para enfrentar
os problemas técnicos, realizar os cálculos e tornar possível o
funcionamento da bomba. Esse domínio científico, porém, não era suficiente
para responder a questões fundamentais: por que a arma ainda deveria ser
construída? Contra quem deveria ser utilizada? A possibilidade de terminar
uma guerra justificaria a morte de milhares de pessoas? Essas perguntas não
pertenciam somente aos campos da Física e da Engenharia, pois envolviam
valores, interesses políticos e decisões éticas. Além disso, embora os
cientistas soubessem que participavam da construção de uma arma, eles não
possuíam controle completo sobre as decisões militares relacionadas ao seu
uso.

Essa situação reforça a posição de Fourez segundo a qual a ciência pode
esclarecer aspectos de um problema, mas não consegue decidir sozinha o que
deve ser feito. No caso da bomba, os cálculos explicavam como construí-la e
permitiam estimar seus efeitos, mas não determinavam se seu uso era
moralmente aceitável. O Projeto Manhattan também exemplifica o caráter
histórico e social da atividade científica: as pesquisas não foram
realizadas de maneira isolada, mas orientadas por um contexto de guerra e
por interesses políticos e militares. A relação com o filme torna concreto
um dos principais méritos do capítulo, que é demonstrar que o rigor de uma
explicação científica não elimina a necessidade de avaliar os valores e os
projetos aos quais ela serve.

O mesmo limite pode ser observado no exemplo de Fourez sobre a afirmação de
que o amor seria apenas uma questão de hormônios. A Biologia pode explicar a
participação hormonal nessa experiência, mas essa explicação particular não
abrange tudo o que o amor representa para os seres humanos. De modo
semelhante, compreender cientificamente o funcionamento da bomba não
significa compreender nem justificar todas as consequências de sua criação.
Nos dois casos, o problema aparece quando uma explicação específica é
apresentada como se respondesse à totalidade da questão. Essa comparação não
diminui a capacidade da ciência de produzir conhecimentos confiáveis; ela
mostra que esses conhecimentos possuem um alcance determinado e não devem
ser tratados como a única fonte de respostas para problemas humanos,
políticos e éticos (FOUREZ, 1995, p. 184--189).

Outro aspecto relevante do capítulo é a maneira como Fourez mostra que uma
mesma descrição pode ser associada a práticas diferentes conforme o meio
social e os valores de quem a interpreta. No exemplo apresentado pelo autor,
a ideia de permitir que outra pessoa disponha de parte de sua intimidade e
criatividade em troca de dinheiro costuma ser relacionada imediatamente à
prostituição, embora também possa descrever um aspecto do contrato de
trabalho. A comparação é provocativa e eficiente porque revela como o efeito
ideológico destaca certas diferenças e deixa semelhanças em segundo plano.
Ela não deve, contudo, ser entendida como se as duas práticas fossem
equivalentes em todos os aspectos. Ao aproximá-las por meio de uma
característica comum, Fourez procura justamente tornar visível algo que as
classificações habituais tendem a ocultar (FOUREZ, 1995, p. 179--180).

Esse exemplo também demonstra por que a discussão do capítulo ultrapassa o
campo científico. Os discursos utilizados no cotidiano orientam a maneira
como diferentes práticas e grupos são julgados, muitas vezes sem que as
pessoas percebam os critérios ideológicos presentes em sua interpretação.
Nesse sentido, uma contribuição importante do texto é incentivar o leitor a
questionar não somente as conclusões apresentadas como verdadeiras, mas
também o ponto de vista, os interesses e as escolhas que permitiram sua
construção.

Embora seja especialmente relevante para estudantes, professores,
cientistas e profissionais que participam da produção ou da aplicação de
conhecimentos técnicos, o capítulo não se limita a um público especializado.
Sua leitura pode ser recomendada a qualquer pessoa interessada em
compreender como as ideias são construídas e legitimadas socialmente. A
principal reflexão deixada pelo texto é que uma sociedade mais crítica não
depende de abandonar a ciência nem de considerar todas as opiniões
igualmente válidas. Ela depende da capacidade de examinar perspectivas
divergentes, identificar seus critérios, confrontá-las com a experiência e
reconhecer os limites da própria visão. Mais do que falar apenas sobre a
ciência, o capítulo defende, assim, uma postura de questionamento e de
responsabilidade diante dos discursos que aceitamos e transmitimos.

% ============================================================
% REFERENCIAS
% Incluir somente obras efetivamente citadas no texto final.
% ============================================================
\section*{Referências}
\begin{singlespace}
\setlength{\parindent}{0pt}

EDITORA UNESP. \textbf{A construção das ciências: introdução à filosofia e
à ética das ciências}. São Paulo: Fundação Editora da Unesp, [s.d.].
Disponível em: \url{https://editoraunesp.com.br/catalogo/9788571390836,a-construcao-das-ciencias}.
Acesso em: 9 set. 2026.

\vspace{\baselineskip}

FOUREZ, Gérard. Ciência e ideologia. \textbf{In}: FOUREZ, Gérard.
\textbf{A construção das ciências: introdução à filosofia e à ética das
ciências}. Tradução de Luiz Paulo Rouanet. São Paulo: Editora da Universidade
Estadual Paulista, 1995. cap. 7, p. 179--193.

\vspace{\baselineskip}

OPPENHEIMER. Direção e roteiro: Christopher Nolan. Produção: Emma Thomas,
Charles Roven e Christopher Nolan. Universal Pictures, 2023. 1 filme
(181 min), son., color.
\end{singlespace}

% ============================================================
% DECLARACAO DE USO DE INTELIGENCIA ARTIFICIAL
% Manter apenas se a disciplina/instituicao solicitar ou permitir.
% Ajustar o texto para refletir com exatidao o uso efetivamente realizado.
% ============================================================
\section*{Declaração de uso de Inteligência Artificial}
\begin{singlespace}
Para a elaboração desta resenha, foi utilizada uma ferramenta de
inteligência artificial como apoio à organização e à revisão textual.
A leitura do capítulo, a seleção das ideias, a análise crítica, a redação
final e a verificação das informações permaneceram sob responsabilidade
do discente.
\end{singlespace}

\end{document}
